\section{Grounded Extractive Memory is Largely Faithful}
\label{sec:faithfulness}

A natural worry is that the LLM extractor writes \emph{hallucinated} facts into
memory that later corrupt answers---a concern formalized for open-ended
generation~\cite{min2023factscore,zhang2023snowball} and, for memory systems,
benchmarked by HaluMem~\cite{chen2025halumem}. We probe whether this is a strong
effect in our grounded setup.

For a sample of extracted facts we ask an LLM judge whether each fact is fully
supported by (stated or unambiguously entailed by) the source session it was
extracted from. In a $120$-fact sample, $0$ were judged unsupported
($\approx\!0\%$). The extractor over-\emph{produces} facts (it inflates the store
and creates the redundancy of \S\ref{sec:redundancy}) but does not
\emph{fabricate}: because it summarizes a provided transcript rather than
generating from parametric memory, it stays close to the source. This contrasts
with the $\sim\!42\%$ unsupported atomic facts reported for open-ended
biography generation~\cite{min2023factscore}.

\paragraph{Takeaway.} For grounded, extractive memory writes, hallucinated-memory
accumulation is a weak effect; the store's problem is \emph{quantity/redundancy}
and downstream \emph{retrieval/reasoning}, not fabrication. (This is a limited
probe---single extractor, single judge, one dataset, conservative
parse-handling---and is not a claim about generative or tool-use memory, where
HaluMem shows hallucination does accumulate.)
